% The ConceptBase User Manual (CB-Manual)
% Copyright 1987-2020 by The ConceptBase Team
% The CB-Manual may only be updated by members of the ConceptBase Team
% as listed on http://conceptbase.cc
% The TeX sources of the CB-Manual and the accompanying
% figures can be freely distributed under the provision that this
% clause is included in file CB-Manual.tex and that CB-Manual.tex
% is the main file of the TeX sources.
% Contact the ConceptBase Team via the URL above if you want to join
% the team.


\chapter{The ConceptBase.cc Usage Environment}
\label{cha:workbench}

The ConceptBase.cc User Interface consists of two main applications:

\begin{description}
\item[CBIva] is the ConceptBase.cc Interface in Java that supports
the editing of Telos frames, displays instances of Telos objects,
etc.
\item[CBGraph] is a graphical editor for Telos models (or modules). Telos models
can be represented by differented graphical types. Insertion and deletion
of Telos objects is also supported.
\end{description}

The interface is based entirely on Java, so it should be usable on all platforms
with a compatible Java runtime environment. The Java interface includes a graphical browser and editor.
Both CBIva and CBGraph can be used as stand-alone Java application.



\section{The workbench CBIva}

CBIva is a textual interface to a CBserver, which emphasizes the use of the frame syntax
of Telos.
You can start the CBIva workbench by the command

\begin{verbatim}
  cbiva
\end{verbatim}

The command will start a script file with the same name.
We assume that
 the installation directory of ConceptBase.cc\footnote{In rare circumstances, you may have to
edit the script files cbiva, cbgraph and cbshell, so that the correct Java Virtual Machine is used and the environment
variable {\tt CB\_HOME} points to the installation directory. See Installation Guide 
in the CB-Forum for details.} is added to the search path
of executable programs. After a few seconds, the CBIva main
window should pop up. CBIva will attempt to connect to a CBserver on localhost or to
a public CBserver (if configured, see \ref{sec:pubcbserver}).
Figure \ref{fig:cbivamain} shows the main window connected
to a CBserver and gives a short description of the buttons in the tool bar.


\cbfigure{cbivamain}{12cm}{Main Window of CBIva}

The main window consists of a menu bar, a tool bar with a button panel, the area for
subwindows and a status bar. The first sub window (Telos editor) contains the history window,
which records all operations for later reuse.
In the following, each component of the user interface is explained.

\subsection{The tool bar}

The tool bar is the button panel below the menu bar. All buttons have tool tips, i.e. small
messages that show the meaning of these buttons. The tool tips appear, if you move your mouse pointer
over the button and do not move your mouse for about one second. The buttons are shortcuts
for some operations that are frequently used and are also available in the menu.
The operations apply to the Telos Editor which has currently the focus.

\begin{itemize}
\item {\bf Quick Connect}: if connected to a CBserver then disconnect; otherwise try to
 connect either to the public CBserver (if configured) or to a running local CBserver; 
 if the connection to a running local CBserver fails, then attempt to start one in the background.
\item {\bf Clear}: Clear the text of the Telos Editor window
\item {\bf Cut, Copy, Paste}: operations on selected portions of the Telos Editor text
\item {\bf Tell, Untell}: The objects specified in the active Telos Editor will be added or removed
from the object base.
\item {\bf Retell}: A new window will popup and ask you to enter the frames that should be untold and
told. The contents of the current Telos Editor will be inserted as default into the two text areas.

\item {\bf Ask}: Evaluate the query specified in the Telos Editor. \footnote{The query can either be a query call
referring to an existing query
or a frame representing a new or existing query definition.} If you specify the name of an ordinary class (i.e.\
not a query class), then ConceptBase.cc will interpret this as a query call to find all instances of that class.
You can also enter an arithmetic expression like\newline
  \centerline{{\tt 100*COUNT(QueryClass)/COUNT(Class)}}
\item {\bf List module}: List the content of the current module as Telos frames
\item {\bf Exit}: Disconnect from the CBserver and exit CBIva.
\end{itemize}



\subsection{The menu bar}
\label{sec:menubar}

The toolbar has button for the most frequent operations of CBIva. The menu bar offers the complete set
of operations including options to open other windows such as CBGraph (see section \ref{sec:graphed}).

\begin{description}
\item[File:]
	\begin{description}
	\item[Connect:] Connect to a ConceptBase.cc server (CBserver) started in another shell/command window (see section \ref{cha:cbserver}).
	\item[Disconnect:] Disconnect from a CBserver.
	\item[Load \& Save Telos Editor:] Load or Save the contents of the current
	Telos Editor.
	\item[Load Model:] Load a source model file (*.sml, *.sml.txt) to the server. As client and server are not supposed
	to share the same file system, this method is now implemented as a normal Tell method. The client
	reads the contents of the file into a string and sends the string to the server.
	\item[Start CBserver:] Opens a dialog and asks for the parameters to start a CBserver.
	If the information has been entered, the server will be started
	and its output will be captured in a separate window. The workbench will
	be automatically connected to that server%
\footnote{This is {\em not\/} the standard way to start the CBserver. Normally, the CBserver
is started in a separate shell/command window as explained in section \ref{cha:cbserver}. The standard
way offers more options and control over the CBserver. See also the installation guide for
a discussion on the various ways to start ConceptBase.}.
	\item[Stop CBserver:] Stops the CBserver, only allowed for the user who started the server
         or who has been designated as administrator user (see option "-a" in section \ref{sec:cbsparams}).
	\item[Close:] Same as Exit if CBIva was not started from CBGraph. If it was started from CBGraph
        and CBGraph is still running, then the Close option only hides the CBIva window.
	\item[Exit:] Exit CBIva, if a CBGraph windows was opened via CBIva, it shall be closed as
        well. Likewise, a CBserver started via CBIva will be closed.
	\end{description}
\item[Edit:]
	\begin{description}
	\item[Clear:] Clears the text area of the currently activated Telos editor.
	\item[Cut,Copy,Paste:] Cut, copy or paste text to/from clipboard.
	\item[Tell,Untell:] (Un-)Tells the text in the currently activated Telos editor
	\item[Retell:] A new window will popup and ask you to enter the frames that should be untold and
told. The contents of the current Telos Editor will be inserted as default into the two text areas.
	\item[Ask Frame:] Temporarily tells the content of the Telos editor, 
        extracts the query names from the frames, asks them
        as query calls without parameters, and returns the result in the Telos editor window
	\item[Ask Query Call:] The query calls%
\footnote{Usually, one only asks a single query call like {\tt Q[v1/p1]}. However, ConceptBase.cc also
supports comma-separated lists of query calls like {\tt Q1[v1/p1],Q2[c1:p2]}. Such lists of query calls
are evaluated one after the other. The results is merged into a single answer. For technical reasons,
calls to builtin query classes like {\tt get\_object[Class/objname]} may only occur as a singleton.}
        (see section \ref{sec:qdefvscalls}) listed in the Telos editor are asked
        and the result is returned in the Telos editor window.
	\item[Load Object:] Load the Telos frame of an object into the Telos editor.
%	\item[View Object as Tree:] Opens a new window and shows the object in tree browser.
	\end{description}
\item[Browse:]
	\begin{description}
	\item[New Telos Editor:] Opens a new Telos editor (see section \ref{sec:telos_ed}).
	\item[Display Instances:] Opens the display instances dialog (see section \ref{sec:display_instances}).
	\item[Frame Browser:] Opens the frame browser window (see section \ref{sec:frame_browser}).
	\item[Display Queries:] Shows the 'visible' (user-defined) queries stored in the current database and provides a facility to
	call these queries (see section \ref{sec:display_queries}).
        \item[Display All Queries:] Like above but includes the many built-in queries of ConceptBase.
        \item[Display Functions:] Shows all functions stored in the current database 
             (see section \ref{sec:display_functions}).
	\item[Query Editor:] Opens the query editor (see section \ref{sec:query_editor}).
	\item[Graph Editor:] Opens the CB Editor (graphical browser, see section \ref{sec:graphed}). If the
	interface is connected to a server, CBGraph will also establish a connection to this
	server and ask for the graphical palette, the initial object to be shown, and the module context (see section \ref{sec:module}). Otherwise, CBGraph will start with no connection.
	\end{description}
\item[Options:]
	\begin{description}
	\item[Set Timeout:] Set the number of milliseconds the user interface waits for a response of
	the server.
	\item[Select Module:] Select the current module (see section \ref{sec:module}).
	\item[Select Version:] Select or create a new version. A version is a special object
	that represents the state of an object base at a specific time. By default, all queries
	are evaluated on the current state of the object base (version "Now"). By selecting another
	version, queries are evaluated wrt. to a previous state of the object base.
	\item[Pre-Parse Telos Frames:] If enabled, the user interface parses the contents of
	a Telos editor before it is send to the server. Thus, syntax errors might be already
	detected at the client side.
	\item[Show Line Numbers:] Enables the display of line numbers in the Telos Editor window.
        Use 'Save Options' to memorize this setting for the next session of CBIva.
	\item[Use Query Result Window:] If enabled, the result of a query is shown in a separate window
	in a table view.
	\item[Look and Feel:] You can switch the look and feel to an other style. Default and preferred is
             'Metal'.
	\item[Save Options:] Saves the current options for subsequent sessions of CBIva.
	\item[Edit Options Manually:] Complete and editable list of options for CBIva. The contents
        is maintained in a file '.CBjavaInterface' of the user home directory.
	\end{description}
\item[Help:]
	\begin{description}
	\item[ConceptBase.cc Manual:] Opens a window\footnote{
Under Linux with Gnome, the default web browser is used to display the content. Otherwise, a Java window capable of
displaying HTML is opened. Due to unknown reasons, the mechanism works for Linux only if at least two browser tools are installed.
You may want to install Chromium as second browser besides Firefox to fulfill the constraint.}
with the online-version of this ConceptBase.cc manual.
	\item[ConceptBase Tutorial I/II:] Opens a window with the online-version of the respective tutorial.
	\item[CB-Forum:] Opens a window to the public version of the ConceptBase Forum with lots of
              examples.
	\item[About:] Shows a dialog with information about this program.
	\item[License:] Displays the license of ConceptBase.cc in a new window.
	\item[CB-Team:] Displays a page about the ConceptBase Team.
	\end{description}
\item[History:]
	\begin{description}
	\item[Load History:] Load previously saved contents of the history window.
	\item[Save History:] Save contents of the history window to a file.
	\item[Redo History:] Redos certain operations which are currently in the history. The operations
	can be selected from a list.
	\item[Set History Options:] Select the type of operations that should be displayed in the history window.
	\end{description}
\end{description}



The preferred way to start a CBserver from CBIva is to use the 'Start CBserver' from the File menu.
It provides an output window for the trace messages of the CBserver, but this window is only visible
if the trace mode is set to 'minimal' or higher. Disconnecting from the CBserver will not stop
it. It has to be stopped explictly or CBIva will stop it if it is shut down itself.
The 'Quick Connect' button of the toolbar provides a simplified way to start a local CBserver.
It does not ask for CBserver parameters. Instead, the default values are used for all parameters
except that the server mode is set to 'slave' and multi-user mode is disabled.
This causes the CBserver to shutdown whenever the last connected client disconnects from it.

You can move the tool bar outside the
main window or display it in vertical form, if you click on the leftmost area of the tool bar
and drag it to another place.




\subsection{The status bar}

The status bar contains some fields that display general information about the status of the application.

\begin{itemize}
\item Connection status, either connected or disconnected from serve, or a pending action
\item Short message, usually about the result of the last action
\item Current version, i.e. the rollback time (see section \ref{sec:representation})
     specified for queries (default: Now)
\item Current module, the database module to which TELL/ASK operations are applied
\end{itemize}

The field for the current version shall display the current time if the
version is set to 'Now'. If set to a rollback time in the past, it displays
both the version name and the time associated to the version. It will
then also hight the background of the text field to make the user aware
that queries are evaluated against a past state of the database.



\subsection{Telos editor}
\label{sec:telos_ed}

The Telos Editor is an editable text area, where you can edit Telos frames.
The operations can be executed from the menu bar or the buttons in the tool bar.
Furthermore, the text area has a popup menu on the right mouse button with
the following items.

\begin{description}
\item[Display Instances:] Displays the instances of the currently marked object
\item[Load Object:] Loads the Telos frame of the currently marked object into the Telos editor
%\item[View Object as Tree:] Opens the tree browser with the currently marked object as root
\item[Display in Graph Editor:] Shows the currently marked object in CBGraph.
\item[Change Module:] Changes the module context to the module path defined by the currently marked text
\item[Clear,Cut,Copy,Paste:] same as in the menu bar
\item[Small,Large:] sets the text font size to either small (12 point) or large (18 point)
\end{description}

If the Telos editor window is empty, then you can drag and drop Telos text files (file
type *.sml or *.sml.txt) into it. CBIva will open the file and paste the content into the Telos
editor window. This function requires that the option 'Show Line Numbers' is enabled
(see section \ref{sec:menubar}). You can also drag and drop URLs pointing to publicly accessible
Telos text files.

\subsubsection{Tell transactions}
The text in the Telos editor window is typically a sequence of Telos frames. When
you press the 'Tell' button, the text is sent to the CBserver and processed there by
the 'TELL' method of the CBserver. This is a single transaction which may fail or succeed.

You can use the string "{\tt \{---\}}" in the text window to indicate that the
frames should be told in multiple transactions. Consider the example:

{\small
\begin{verbatim}
Shop in Class end
Guest in Class with
      attribute
            dept: Shop
end
{---}
GuestEmployee in Class isA Guest,Employee end
\end{verbatim}
}


Activating the 'Tell' button instructs CBIva to split the text into two parts and tell them in two
separate transactions to the CBserver. This feature is useful when some parts contain meta formulas
(see section \ref{sec:CCmeta})
that first need to be compiled before they are used in subsequent parts of the whole text.
If you do not use meta formulas, then you can omit this feature.


\subsection{History window}
\label{sec:history_window}

The history window is part of the main Telos editor. It stores all operations and their
results, so that they can later be used again. The buttons scroll the history back or forward,
copy the text into the Telos editor or redo the operation in the history window (see figure \ref{fig:historybuttons}).
If the current operation is an ``ASK'', then a single click on the copy-button will copy the
query to the Telos Editor, and a double click will copy the result of the query.
The size of the history window can be reduced by using the slider bar between the
Telos editor and the history window.

\cbfigure{historybuttons}{0.6\textwidth}{Buttons of the History Window}

\subsection{Display instances}
\label{sec:display_instances}

This dialog displays the instances of a class. The class is entered in the
text field. When you hit return or press the ``OK'' button, the instances of this
class will displayed in the listbox.

If you double click on an item in the listbox, the instances of this item
will be displayed. The frame of a selected item can be loaded into the Telos
editor by clicking the ``Telos Editor'' button.
A history of already displayed classes is stored in the upper right
selection list box. The ``Cancel'' button closes the dialog.

\subsection{Frame browser}
\label{sec:frame_browser}

The Frame Browser (see figure \ref{fig:framebrowser}) shows all information relevant to one object in one
window. The window contains several subwindows with list boxes that show
super- and subclasses, the classes, the instances, attributes and
objects refering this object. In the center of the window, a small
window with the object itself is shown.
To view the attributes of the object, you must first select the
attribute category in the subwindow ``Attribute Classes''.

\cbfigure{framebrowser}{11cm}{Frame Browser}


The Frame Browser can be used with and without a connection to
a CBserver. If it is not connected, it retrieves the
information out of a local cache, which can be loaded from a file
by using the ``Load'' button. The file has to be plain text file
with Telos frames. All objects in the cache can be saved into a text file
as Telos frames with the ``Save''
button. The contents of the cache can be viewed
with the ``Cache'' button. The result of a query can be added to
the cache by using the ``Add query result'' button.

The button ``Telos Editor'' inserts the Telos frame of the current
object into the Telos Editor.

\subsection{Display queries}
\label{sec:display_queries}

This dialog displays all visible queries%
\footnote{Visible queries are those queries that are not instantiated to the class
{\tt HiddenObject}. Functions and certain system queries are excluded from the display.}
stored in the current object base (see figure \ref{fig:displayqueries}). 
From the list box, you can
select a query and ``ask'' it or load its definition into the Telos editor.

\cbfigure{displayqueries}{11cm}{Display queries dialog}

If you ``ask'' a generic query class with parameter, another dialog will
ask you to specify the parameters. For each parameter, you can specify
whether the value entered should be used as ``substitute'' for the parameter
or as a ``specialization'' of the parameter class (see section \ref{sec:CBQL}).
You can select a value for the parameter from the drop-down list if you have
clicked on the ``Show Values'' button. Note that this list might be very long.
Especially for the predefined queries it usually returns all objects in the database
as any object can be used as a parameter for these queries.

\subsection{Display functions}
\label{sec:display_functions}

This dialog is similar to the previous dialog but displays instances
of {\tt Function}. Note that functions are formally special queries.
Consequently, the dialogs for functions and queries are pretty much the same.
The separation into two dialogs serves quicker handling.


\subsection{Query editor}
\label{sec:query_editor}

The Query Editor (see figure \ref{fig:queryeditor}) allows the interactive definition of queries.
The name of the query is entered in the upper left text field,
the super class in the upper right field. After you have entered
this information, the list box ``Retrieved Attributes'' will be
filled with all available attributes.

Now, you can select the attributes you want to have in the result.
For selection of more than one attribute, you must press the
CTRL key and select the attribute with a mouse click at the same time.
All attributes can be deselected by the popup menu.

\cbfigure{queryeditor}{11cm}{Query Editor}

In the right listbox you can add computed attributes. The right
mouse button brings up a popup menu, which lets you add or delete
an attribute.

In the text area below the two list boxes, you can add
a constraint in the usual CBQL syntax. The constraint must
be enclosed in \$ signs.

The text area below, shows the Telos definition of the query
and is updated after every change you have made. If your query is
finished, you can press the ``Ask query'' to test the query,
i.e. it is told temporarely and the results are shown in separate
window. If you are satisfied with the result you can press the Tell
button to store the query in the object base.

\subsection{Tree browser}
\label{sec:tree_browser}

{\bf Note: The tree browser has been removed from CBIva due to legacy dependencies on
packages no longer supported by Java.}

The Tree browser (see figure \ref{fig:frametree}) displays the super classes, classes and
attributes of an object in a tree. To start the tree browser,
you must mark an object in the Telos editor and select the
item ``View Object as Tree'' from the popup menu or the Edit
menu.

\cbfigure{frametree}{9cm}{Tree Browser}

To expand an item, just double click on the icon. If you mark an object
name in the tree, you can load into the Telos editor with the
``Telos Editor'' button or open a new tree browser with the
button ``View Object as Tree''.


\newpage
\section{The graph editor CBGraph}
\label{sec:graphed}

The CBGraph Editor is an advanced graphical modelling tool that
supports the browsing and editing of Telos models.
It supports user-definable graphical types, i.e.\
objects may be visualized by dedicated graphical layouts. In
addition to predefined graphical types, the user can add his/her own
graphical types by modifying and adding certain objects in the knowledge
base. Furthermore, the standard components can be replaced by own classes
implementing specific application-dependent behaviour.

In the following, we first give an overview of CBGraph application
and then present the main components and functions of CBGraph.
Details about the use of graphical types can be found in Appendix
\ref{cha:graph-typen}. An example for the definition of graphical types
of the Entity-Relationship model is given in Appendix \ref{sec:ER-diagrams}.

\subsection{Overview}

The CBGraph Editor is entirely written in Java. It can
therefore be used on any platform with Java 1.4 or compatible successors of Java 1.4.
CBGraph is integrated with CBIva, i.e.\ Telos frames
of objects shown in CBGraph can be loaded directly
into a Telos editor and vice versa. CBGraph allows to open
several 'internal windows' in its main window. Each internal window has a separate
connection to a CBserver. Thus, within a CBGraph
you can establish multiple connections to the same CBserver or even to
different servers.

The communication with the CBserver is done using pre-defined
ConceptBase.cc queries and a special XML-based answer format.
CBGraph requests information about the objects (names, attributes, etc.)
but also about their graphical type.

\cbfigure{overview}{0.9\textwidth}{CBGraph with three internal windows}

Figure \ref{fig:overview} gives an
overview of the CBGraph Editor. Three internal windows
have been opened, the two windows on the left have been connected
with the same server running on ``localhost'', port 4001. The small
window in the upper right corner is not connected to a server, but
a few objects have been created. The title of the internal windows 
displays the graphical palette and the current module (if connected
to a ConceptBase server) plus the connection status (either 'offline'
or the hostname and portnumber of the ConceptBase server). The content of
an internal window is a graphical view on the database module of the
ConceptBase server it is connected to. It is possible that different
internal windows connect to different ConceptBase servers, though this
is not a typical use of CBGraph. If you start CBGraph from CBIva, then
you can only start a single instance of it. However, you can start any
number of CBGraph instances via the 'cbgraph' command (see below).


The two left internal windows of figure \ref{fig:overview} are connected to the same
server and show the same model but in different representations. This is caused by the fact,
that for the upper window, the default graphical palette has been
chosen, and for the lower window, a customized graphical palette specifically
designed for the ER model has been selected (see appendix \ref{sec:ER-diagrams}
and \ref{cha:graph-typen}).

Furthermore, you can notice that the object ``QueryClass'' is represented
by two different components. In the upper view, the detailed {\em component view}
has been activated by a double-click on the object. It shows the Telos
frame of the object. By a double-click on the title bar of this component,
one can switch back to the default view of this object. Thus, each object
can be shown by a small component (the default view) and a large
component which gives more detailed information. Components are in this
context specific Java objects, namely instances of javax.swing.JComponent.
Thus, different components can be provided to represent an object (e.g.,
tables, buttons, text fields). You can implement your own component
and integrate into CBGraph by extending a specific Java class.
Details about the customization of CBGraph using graphical
types and other components can be found in appendix \ref{cha:graph-typen}.

CBGraph can be used to edit Telos models (see section
\ref{sec:graph_edit}). It can also display implicit
relationships between objects, which have been derived by rules or Telos
axioms. For each type of relationship (instanstation, specializations, and
attributes), one can choose to see only the explicit relationships or to see
all relationships.

\subsection{Starting CBGraph via CBIva}

CBGraph can be invoked from CBIva
via the menu item {\em Browse} $\rightarrow$ {\em Graph Editor}.
If CBIva is connected with a CBserver, CBGraph will be
connected to the same server and you will be prompted to enter the object
name you want to start with, the name of a graphical palette, and the database module. The
graphical palette is a Telos object which represents a set of
graphical types which will be used to visualize Telos objects (see
Appendix \ref{cha:graph-typen}). On startup, CBGraph
retrieves all information about the graphical types from the CBserver.
If CBIva has no connection with
a server, CBGraph will be started without a connection
and no internal window will be opened within CBGraph.

The connection to a CBserver can also be established via the File / Connect
menu. The dialog box has two tabs. The first is for providing the host
and port number number of the CBserver. The second is for providing the 
start object (default "Class") to be displayed, the graphical palette (pick from
a list), and the database module (default "oHome").

The most comfortable way to start CBGraph is to double-click the graph file. It is equivalent
to starting it via the command 

{\small
\begin{verbatim}
  cbgraph graph.gel
\end{verbatim}
}

To do so, you need to configure your desktop according to the instructions at
\url{http://conceptbase.sourceforge.net/CB-Mime.html}.



\subsection{The {\tt cbgraph} command}
\label{sec:cbgraphcmd}
You can also start CBGraph as a stand-alone utility. The command

{\small
\begin{verbatim}
  cbgraph [options]
\end{verbatim}
}

will start an unconnected graph editor. You can interactively connect it to a running CBserver
and open new windows to display graphs. More interesting is the use with a stored graph file,
or several graph files, resp.:

{\small
\begin{verbatim}
  cbgraph [options] filename [filename ...]
\end{verbatim}
}

The format of the graph file is called Graph Editor Layout (GEL). It stores not only the layout of nodes
and links but all other data necessary to edit the graph objects. In particular, it contains
connection details of the CBserver module from which the graph was created. You can open more than one
GEL file. Each will be loaded in its own frame. 

The are a few options for the 'cbgraph' command synchronizing the data stored in the graph file
with the CBserver:
\begin{description}
\item[+r] With this option CBGraph will be instructed to load the module sources on its current module
  path from the CBserver and save it to the graph file when the 'save' function is enacted. 
  The 'System'  module will not be saved since it is (typically) not changed. For example, if CBGraph 
  displays objects of module "System-oHome-M1", then the save function will store the module sources
  of 'oHome'  and 'M1'  to the graph file. We say "CBGraph reads the sources from the CBserver and
  saves them to the graph file". This option disables writing the module sources from the graph file to the 
  CBserver when CBGraph is started. An exception is applied when CBGraph needs to start a local 
  CBserver on the fly. Such a CBserver has an empty database (except system objects) and
  the module sources are told to such a fresh CBserver as if the "+rw" flag had been applied.
\item[+w] This option enables the reverse direction. It will instruct CBGraph to extract module sources
  from the given graph file (or any graph file that is loaded via CBGraph) and tell them to the CBserver
  that it is connected to. You have to start a CBserver as a separate process using the hostname and
  port number that is stored in the graph file. This option disables by default including the module
  sources to the graph file when saving the graph file.
  The CBserver can well have an empty database because all required user-defined objects are stored
  as sources in the graph file. If the database is not empty, the definitions from the sources are
  added to the current database. It may well be that the operation causes an integrity violation or
  fails because the current user has no write permission on a given module. 
  In case of success we say "CBGraph loads the module sources from the graph file and writes
  them to the CBserver".
\item[+rw] Enable bidirectional synchronization  (default).
\item[-rw] No synchronization.
\item[+f]  Enable bidirectional synchronization and also write the module sources to text files in the directory
   where your started CBGraph; useful for debugging.
\item[-demo] Disables certain menu items such as File/Save to make CBGraph more stable in demonstration
   scenarios.

\end{description}

If you supply these options when calling CBGraph, it will also store them in the graph file when you store it.
If you later call CBGraph with such a graph file, it will apply the stored options unless you specify
new options in the command line. Hence, the options stored in the GEL file serve as new defaults for
synchronizing the graph file with the CBserver.

Consider the two calls of CBGraph below. The first call sets the synchronization option to "+r", i.e.\
stores the module sources from the CBserver in the graph file when it is saved. 
The second call has no such option. Hence, it shall adopt the "+r" that was stored earlier in the graph file
{\tt example.gel}.

{\small
\begin{verbatim}
  cbgraph +r example.gel
  cbgraph example.gel
\end{verbatim}
}

This behavior is useful if you keep the models in a persistent CBserver and want to use the module sources in
the graph file only as a backup storage. It will then not tell the module sources from the graph file to the CBserver.
The TELL operation can be very costly and is redundant if the CBserver anyway has all definitions already stored.



If you provide a filename (or several), then it must have been previously created
by another graph editor, e.g.\ a graph editor started via CBIva.
The above command will then display the graph stored 
in the graph file in an internal window and attempt to connect to the same
CBserver that was active when the graph file
was created. Hence, the graph file is a materialized view on the CBserver database that is 
visualized with 'cbgraph'. You edit a graph file with 'cbgraph' like you are editing a drawing
with a drawing tool. The only difference is that the graph is linked to the database.

If the CBserver module specified in the graph file is not accessible, the graph is still opened
and you can edit it. You can then however not commit changes to the database or add new
objects from the database. The CBGraph editor displays the connection status in the title
of its internal window containing the graph.
The hostname of a CBserver in a graph file is by default 'localhost'. Consequently, the graph
editor will try to establish a connection to a CBserver running on localhost. If you want to connect
to a remote CBserver, then you should specify in CBIva the full domain name of the CBserver,
e.g. 'myhost.acme.com'. This long name will then be stored in the graph file that is created
from a graph editor started via CBIva. Such graph files can then be copied to other computers
and can be loaded with CBGraph to auto-connect to the remote CBserver specified in the graph file,
provided that CBserver is running and accessible.

If no CBserver is accessible, CBGraph will attempt to start a local CBserver in the background
provided that the graph file specifies "localhost" as hostname and the graph file contains module sources.
In other cases, CBGraph switches to the offline mode. You can still change location
and size of the graphical elements and store it back to the graph file. But you cannot delete
objects and you cannot add objects to the database. The menus to show attributes, instances,
and subclasses shall also not work.

An example on how to create and use graph files containing materialized database views
is available in the CB-Forum at
\url{http://merkur.informatik.rwth-aachen.de/pub/bscw.cgi/3613919}.

\subsection{Redirecting the CBserver location}

You can use the command line argument '-host' to override the hostname and portnumber encoded in
the graph file. Assume a graph file 'graph1.gel' was created from a a CBserver connection
at 'localhost:4001'.
Then, loading this graph file in a subsequent call will also connect to 'localhost:4001'. If you
want to use another CBserver, e.g.\ running on 'myhost.acme.com:4002', then call CBGraph 
from the command line as follows:

{\small
\begin{verbatim}
  cbgraph +rw -host myhost.acme.com:4002 graph1.gel
\end{verbatim}
}

Note that the CBserver must be running at the remote location before you enter the above command.
When you subsequently save the graph file, it will have 'myhost.acme.com:4002' encoded as its
connection. The redirection also works in the reverse direction. So, assume that the graph file
'graph2.gel'  was created for a connection at 'myhost.acme.com:4002'. Then, the following command
will redirect the connection to localhost:

{\small
\begin{verbatim}
  cbgraph +rw -host localhost graph2.gel
\end{verbatim}
}

The default port number is 4001. If no CBserver is running at 'localhost:4001', then CBGraph shall start
it in the background. Note that local CBservers can currently only be started on Linux hosts.



\subsection{Moving objects}
\label{sec:moveobj}

All explicit information is a proposition in ConceptBase (see section \ref{sec:representation}),
including attributes,
specializations and instantiations. You can move objects in the graph 
editor by pointing the cursor to the object's label and then dragging it
while keeping the left mouse button pressed. 
If a graph has many nodes and edges, then it is recommended to first click
once on the node or edge to be dragged and then to drag it. This will
switch off the anti-aliasing while dragging and thus be faster. 

You may also want to select a group of objects and then move it as a
whole. To do so either span a selection box with the left mouse button around the
object to be selected, or press the "shift" key and select multiple objects
individually.
After a move, CBGraph shall redraw dependent edges that might be misplaced.

Some edges like instantiation links have no label in CBGraph,
depending on the graphical type associated to it (see appendix
\ref{cha:graph-typen}). In such cases, a small square dot is displayed 
on the edge. Click on this square dot to select the edge and to drag it.
If the edge has a background color (see appendix \ref{cha:graph-typen}) different from
the edge color, then the square dot is drawn in the background color. 

Moving nodes and edges can sometimes lead to quite messy edge curves where
the edge middle point is distant from the middle of the straight line
between the source and destination of the edge. You can clean up such graphs
by pressing the "shift" and "control" keys together and then click on a
node whose edges shall be straightened. 







\subsection{Menu bar}

The menu bar provides access to the most important functions
of CBGraph.

\subsubsection{File menu}

\begin{description}
\item[Connect to server:] Connect to a new server. You can have
multiple connections to one server or different servers at the same
time. Each connection will be represented in one internal window
with the graphical palette, database module, host name and port number
in the title bar of the window.
The connection dialog consists of two tabbed panes. In the first one
(Address), you can enter the host address (name or IP number) and the
port number. In the second pane (Initial Object), you can specify
the initial object to start the browsing process, the graphical
palette, and the database module.

\item[Start CBIva Workbench:] Start a CB workbench (aka CBIva). If you
started CBGraph directly or if you have already closed
the workbench window, you can (re-)start the workbench by this menu
item. If connected to a CBserver, the new CBIva will list the content
of the current module in its Telos editor window.

\item[Save:] Save the current graph into a file. The current
nodes, their location and the links will be saved into a file which
can be reloaded later. The graphical types will also be saved into the file. 
By default, the file will get the extension ``gel'' (Graph Editor Layout).
A checkbox lets to select whether you want to save the module sources
into the graph file. It overrides the '+/-rw' options of the cbgraph command.

\item[Load:] Load a graph that has been saved with the previous menu
item. Existing nodes and links in the current window will be erased.
If the graph file was previously created to include module sources,
then the module sources are told to the CBserver before the graph is displayed\footnote{
If you load a graph file that was created with different graphical types
than the ones defined in the CBserver module that CBGraph might be connected to,
then the graph file is inconsistent with the current CBserver module.
You may be able to repair the inconsistency via the menu option
Current connection / change graphical palette. Likewise, some objects displayed
in the graph could be undefined in the database. You can validate them by one
of the validation tabs in the menu "Current connection". }.
A checkbox lets to select whether you want to load the module sources from
the graph file and tell them to the CBserver. 

\item[Print:] Print the current graph. If the graph is larger than
the page size, it will be automatically reduced to fit into the page.
Thus, all printouts will be on one page.

\item[Save image of graph:] Saves an image of the current graph
as PNG or JPG file. The PNG file format should be preferred as it
delivers better results. For technical reasons, the image saving has
to set the background color of the node/edge labels. This can overlap
with the node shapes when they are not rectangles. If so, then use the
screenshot function of your operating system to get better results.

\item[Close:] Closes CBGraph. A CBIva window will be still available
if it has not been closed before.

\item[Exit:] Exit CBGraph and CBIva. This operation will close
both windows of the ConceptBase.cc User Interface (CBIva), if they were started
in combination, and exit the program.
\end{description}

If the graph was loaded from a GEL file and it was edited in the session,
CBGraph shall ask the user whether to save the edited graph to the file
when the user terminates CBGraph.



\subsubsection{Edit menu}

The operations in this menu have an effect on the selected objects.
You can select an object by clicking on it with the left mouse button.
Multiple selection is possible dragging a rectangular area while holding
down the left mouse button or by holding the Shift-key and clicking
on objects.

\begin{description}
\item[Erase Selected:] This option will remove the currently selected
nodes and edges from the view. This operation has no effect on the database.
\item[Selection:] With this submenu, you can either select all objects,
all nodes or all edges in the current frame. Furthermore, you can clear your
current selection.
\end{description}

\subsubsection{Options menu}
\label{sec:optionsmenu}

The options will be stored in a configuration file (see section \ref{sec:config})
when you exit CBGraph.
\begin{description}
\item[Language:] The text for menu items and buttons is available
in two languages (German and English). With this option you can switch
between the languages. This option is currently without function. All
menu labels are set to English.
\item[Background Color:] Here you can change the background color of
the graph to your favorite color.
\item[Component View:] With this option, you can configure the view
of an object if the component view is activated by a double click or
by the popup menu. By default, a tree-like representation of the
object with its classes, instances, super classes, subclasses, outgoing
and incoming attributes is used. If you select ``Frame'', then the
Telos frame of the object will be shown in a text area, (see figure \ref{fig:componentview}).
\item[Invalid Telos Objects:] CBGraph can validate objects
that are currently shown in the graph, i.e.\ it checks whether the object
is still valid in the database or it is has already been removed (see Current
Connection menu). If you select ``Mark'' here, then the objects will be marked
with a red cross as invalid. ``Remove from display instantly'' will remove
the objects directly from the view.
\item[Popup Menu:] These options control the behaviour of the Popup Menu
(see section \ref{sec:graphed:PopupMenu}). The delay is the time (in milliseconds)
an item of the popup menu has to be selected before the submenu is shown. Note that 
the construction of a submenu might require a query to the CBserver.
If the option {\em ``Popup Menu blocks while waiting for server''} is activated, then
the editor will block the UI while it waits for an answer of the CBserver.
Otherwise, the query to the server will be executed in a separate thread, and
interaction with the UI will be possible. If you have the problem that some
submenus are still shown after you have used the popup menu, set the delay to
0 and activate the option {\em ``Popup Menu blocks while waiting for server''}.
\item[Look \& Feel:] This option allows you to adapt the look and feel of CBGraph
windows to your desktop environment. By default, the 'Metal' look and feel is enabled.
CBGraph is not yet compatible with all Java look and feels. We thus recommend to
stick to the default 'Metal' look and feel.
\item[Enable Click Actions:] Click actions automatically trigger an active rule
(see section \ref{sec:eca}) when nodes with a certain graph type are pressed
(see section \ref{sec:clickactions}).
You can enable and disable this feature. If you want to move nodes inside a
graph that has nodes with click actions, then you may want to disable click actions
until you have fnished the re-arrangement of nodes.
\item[Enable Derived Links:] By default, CBGraph can display derived links/relations
to and from a given object. If the database is large and the rules are complex, 
this feature becomes almost unusable due to long delays. In such case, disable the
display of derived links. The behavior of CBGraph is then different wrt.\ the
pop menu displayed when right-clicking on a node or link. It would show all possible
attribute categories rather than only the used ones. For technical reasons, a change
in this flag becomes only effective after re-starting CBGraph.

\end{description}

\cbfigure{componentview}{9cm}{Component view of Telos objects: tree and frame}

If an object is displayed in component view, one can switch back to the node
view by double-clicking its title section. 

\subsubsection{View menu}

CBGraph has an experimental layout algorithm that may be useful
to reorganize the layout of a complex graph. The heuristic of the layout
algorithm is rather simple and does not minimize link crossings.

\begin{description}
\item[Enable automatic layout:] Call a layout algorithm everytime the graph
is changed, e.g. by expanding the attributes. Disabled by default. Enabling
it is not recommended.
\item[Undo last layout operation:] Undo the last  change of the
graphical layout. This option is only available when automatic layout is enabled.
\item[Layout graph:] Call the layout algorithm.
\item[Zoom:] Set the zoom factor of the graph, e.g. 120 (20 percent enlarged).
\item[200,100,75,50,25:]: Set the zoom factor accordingly.
\end{description}


\subsubsection{Current connection menu}

These operations have effect on the current connection, i.e. the connection
of currently activated internal window.

\begin{description}
\item[Query to server:] This operation will open a dialog which prompts
you to enter the name of a query (see figure \ref{fig:querytoserver}).
The query can also be parameterized.
If you click on the ``Submit Query'' button, the query will be sent to
the server and the result will be displayed in the listbox. You can select
the objects that should be added to the graph. Selection of multiple objects is
possible.
\item[Validate and update shown objects:] This operation will check for
every object, if it is still valid (i.e. if it still exists in the database),
update the graphical type of the object, and the internal cache of the
object is deleted (see below, section \ref{sec:graphed_cache}). Depending on the option
``Invalid Telos Objects'' (see above), the objects will be either marked or removed
from the current view.
\item[Validate and update selected objects:] Similar to the previous option but
is only applied to objects in the graphical view that have been selected.
\item[Change graphical palette:] The current graphical palette (=assignment of
graphical types to nodes and links) is replaced by another one. 
\item[Change module:] Change the database module for the currently active internal window.
\end{description}

\cbfigure{querytoserver}{6cm}{Query dialog}


\subsection{Tool bar}

The tool bar (see figure \ref{fig:graphed_toolbar})
consists of a set of buttons that are mainly short cuts
for some menu items. The right half of the tool bar provides buttons
for the creation of Telos objects.

\cbfigure{graphed_toolbar}{\textwidth}{Tool Bar of CBGraph}

\begin{description}
\item[Open a new frame (without connection):] Opens a new internal window without a connection
to a server. Within this internal window, you can create new Telos objects, load existing
layouts and save new layouts. Information about the graphical types is loaded
from an XML file included in the JAR file\\ (\verb+$CB_HOME/lib/classes/cb.jar+).
\item[Load graph:] see {\em File Menu} $\rightarrow$ {\em Load}
\item[Save graph:] see {\em File Menu} $\rightarrow$ {\em Save}
\item[Hide selected objects:] Hides the selected objects from the current
view. This operation has no effect on the current database, i.e.\ the objects
will be not deleted from the database.
\item[Open connection:] Opens a new internal window with a new connection to
a server. See {\em File menu} $\rightarrow$ {\em Connect to server}.
\item[Show object:] This operation adds an object defined in the database to the graph.
You will be prompted to enter the object name of the object you want to
add to the graph.
\item[Show links between marked objects:] This operation will search for
relationships between the selected objects. Do not select too many objects
for this operation, $n^2$ queries have to be evaluated for $n$ objects.
\item[Creation of objects:] The following four buttons open the ``Create Object''
dialog to create new individual objects, attributes, instantiations, and specializations.
See section \ref{sec:graph_edit} for more details.
\item[Show added/removed objects:] Shows the objects that have been added or
removed since the last commit (or since the window has been opened).
Here, you can also select objects to undo the change, i.e. remove added objects
or re-insert removed objects.
\item[Commit:] Sends the changes to the server. The list of objects to be added
or removed is transformed into a set of Telos frames and transferred to the
server. This button is highlighted when some change (added or removed object)
can be committed.
\end{description}

\subsection{Popup menu}
\label{sec:graphed:PopupMenu}

The popup menu is activated by a click on the right mouse while the
cursor is located over an object.

\begin{itemize}
\item {\bf Toggle component view:} \\
	switches the view of this object. In the detailed component view,
	you can either see the frame of this object or tree-like
	representation of super- and subclasses, instances, classes, and
	attributes of this object (see figure \ref{fig:componentview}).
\item {\bf Display in Telos Editor:}\\
    This operation will load the frame of the object into the
    Telos editor. If the Telos editor contains already some text, then
    the frame is appended to the existing text.
\item {\bf Super classes, sub classes, classes, instances:}\\
    for each menu item you can select whether you want to
    see only the explicitly defined super classes (or sub classes, etc.)
    or all super classes including all implicit relationships.
    The query to the CBserver to retrieve this information
    will be done when you select the menu item. So, the construction
    of the corresponding submenu might take a few seconds.
\item {\bf Incoming and outgoing attributes:}\\
    CBGraph will ask the CBserver for the attribute
    classes that apply to this object. For each attribute class, it is
    possible to display only explicit attributes or all attributes as above.
    The attribute class ``Attribute'' applies for every object and all
    attributes are in this class. Therefore, all explicit attributes of
    an object will be visible in this category. However, there will be
    no attributes shown in the ``All'' submenu, as it would take too
    much time to compute the extension of all implicit attributes.
\item {\bf Add Instance, Class, SuperClass, SubClass, Attribute, Individual:}	\\
    These menu items will open the ``Create Object'' dialog where
    you can specify new objects that should be created in the database.
    Note, that these modifications are not performed directly on the
    database. The editor will collect all modifications and
    send them to the CBserver when you click on the ``Commit''
    button. See section \ref{sec:graph_edit} for more details.
\item {\bf Delete object from database:} \\
    This operation will delete the object from the database. As for the
    insertion of objects before, the modification will be send to
    the server when you click on the ``Commit''
    button. Note that this operation has an effect on the database
    in contrast to the next operation.
\item {\bf Hide object from view:} \\
    The object will be removed from the current view. This operation
    has no effect on the database, i.e.\ the object will not be deleted
    from the database.
\item {\bf Show in new Frame:}\\
    A new internal window (within CBGraph) will be shown
    and the selected object will be shown in the new window.
\item {\bf Straighten attached edges:}\\
    All edges starting from and ending in the selected node are made straight and the edge labels move to the middle
    of the edge. This is the same function as described in section \ref{sec:moveobj}.
\item {\bf Freeze / Unfreeze:}\\
    The position of the selected object is frozen (if not yet frozen). A "frozen" object cannot be moved or double-clicked.
    If frozen, the selected object is "unfrozen".   This function is useful when certain objects should not be moved,
    e.g.\ when they are serving as regions, in which other objects are positioned.
\end{itemize}


\subsection{Editing of Telos objects}
\label{sec:graph_edit}

CBGraph supports also the creation and deletion of Telos
objects. A Telos object in the context of CBGraph
is a {\em proposition} as described in chapter \ref{cap:language}.
As described there, there are four types of Telos objects:

\begin{itemize}
\item Individuals,
\item Instantiations (InstanceOf),
\item Specializations (IsA), and
\item Attributes.
\end{itemize}

\cbfigure{createobject}{0.8\textwidth}{Create Object dialogs for Individuals, Attributes, Instantiations, and Specializations}

For each object type, we provide a dialog to create this object
type as shown in figure \ref{fig:createobject}. This dialog is opened
by clicking on one of the ``Create'' buttons in the tool bar, or by selecting
an ``Add ...'' item from the popup menu (e.g, ``Add instance'' or ``Add subclass'').
If there are some objects selected in the current internal window, then the object names of
these objects will be inserted into the text fields of the dialog in the order they
have been selected (i.e., the first text field will contain the name of the object which
has been selected first). Furthermore, if you move the cursor into a text field in the ``Create Object''
dialog and select an object in the graph, then the name of this object will be inserted
into the text field.

As changes might lead to a temporary inconsistent state of the database, we
do not execute the changes directly on the database. They are stored in an internal
buffer in CBGraph and executed when you hit the ``Commit'' button in
the tool bar.

\begin{description}
\item[Creating Individuals:] If you want to create a new individual object,
  you just have to specify the object name. You have to enter a valid
  Telos object name, for example it must not contain spaces. In addition, you can select a
  graphical type for the object. Note that the selection of the graphical type
  has no effect in the database, e.g. by selecting the graphical type of a class ({\tt ClassGT})
  the object will not declared as an instance of {\tt Class}. If you have performed
  the commit operation, the object will get the ``correct'' graphical type from the server.

\item[Creating Instantiations:] In the dialog for instantiations, you have to
enter the name of the instance and the name of the class in the two text fields.
If the object entered does not yet exist, it will be created and represented in the default
graphical type.

\item[Creating Specializations:] This dialog is similar to the one before except that
you specify here the name of the subclass and the name of the superclass. Objects that
do not exist yet, will be created and represented in the default graphical type.

\item[Creating Attributes:] This is the most complex dialog as you have to specify
the source, the label, the value, and the category of the attribute.
The source and the value (or destination) of the attribute are normal object names.
The label may be any valid Telos label. The attribute category has to be a select expression
specifying an attribute category (e.g., {\tt Employee!salary}, see chapter \ref{cap:language}).
The attribute category can be selected from a listbox by clicking on the ``Select'' button
next to the text field of the attribute category. All attribute categories that apply to
the current source of the attribute will be shown. Note that the list will be empty
if the source object does not yet exist in the database.

If you have specified the attribute category, you can also select the attribute value
from a listbox by clicking on the ``Select'' button next to the text field for the attribute
value. The listbox will show all instances of the destination of the attribute category
(e.g. all instances of {\tt Department} for the category {\tt Employee!dept}).

If you select the radio button ``Show Attribute Instantiation'' then CBGraph
will also show the instantiation link for the attribute. For example, if you create
a new attribute for {\tt John} with the label {\tt JohnsDept} in the attribute category
{\tt Employee!dept}, then the instantiation link between {\tt John!JohnsDept} and
{\tt Employee!dept} will also be shown. As the graph gets quite confusing with too
many links, this radio button is not selected by default.
\end{description}

Deletion of objects is also possible. As this operation should not be mixed up
with the removal of an object from the current view, this operation is just
available from the popup menu (item ``Delete Object from Database'').

As you might make mistakes while editing the model, there is the possibility
of undoing changes. The button ``Show added/removed objects'' list all objects
that have been added or removed (since the last succesful commit or since the connection has been
established). A screenshot of the dialog is shown in figure \ref{fig:graphed_addremove}.
The left list shows the objects that have been added, the right list shows the objects
that have been removed. By clicking on the button ``Re-Insert/Delete'' object, the selected
objects will be re-inserted in or deleted from the graph \footnote{You can unselect an object
by holding down the Control key and clicking on the object.}.

\cbfigure{graphed_addremove}{7cm}{List of objects which have been added or removed}

If you are satisfied with the changes you have done, you can click on the
``Commit'' button. Then, CBGraph will transform the objects to be
added or removed into a list of Telos frames and send them to the server
using the TELL, UNTELL, or RETELL operation. If the operation was successful,
all {\em explicit} objects will be checked if they are still valid and if their graphical
type has changed (as in the ``Validate and update'' operation from the ``Current Connection''
menu). If there is an error, the error messages of the server will be displayed
in a message box. The internal buffer with the objects to add or remove will be
not changed in this case.

\subsection{Caching of query results within CBGraph}
\label{sec:graphed_cache}

To improve the performance of CBGraph, several caches
are used. On the other hand, the use of a cache causes
several problems which will be addressed in this section.
In particular, the caches of CBGraph are not
updated automatically if the corresponding data in the
server is updated.

\begin{description}
\item[Graphical Palette and Graphical Types:]
When a connection to a server is established, CBGraph
loads the graphical palette and all its graphical types including
their properties and other information. If an
object has to be shown, the server sends only the name of the graphical
type, the information about the properties are taken from
the cache. Thus, if you change the graphical palette
or a graphical type after CBGraph has established the
connection, this change will not be visible in CBGraph.
There is currently no method implemented to update the cache manually.

\item[Graphical Types of Objects:]
When an object is loaded from the server also the graphical type
for this object is retrieved. The graphical type of the object
is updated when you invoke the ``Validate and Update'' operation from the
``Current Connection'' menu.

\item[Lists of super/sub classes, classes/instances, attributes:]
The lists in the popup menu or in the tree-like view of an object
are produced by evaluating queries. To reduce the communication between
client and server, each query will only be evaluated once (when the corresponding
popup menu should be shown or when the part of the tree should be shown).
The result will be stored in a cache for each object. This cache is emptied
when you invoke the ``Validate and Update'' operation from the
``Current Connection'' menu.
\end{description}


\subsection{Graph files}
\label{sec:gel}

Graph files (extension 'gel') are binary files that store the current state of
a graph displayed in CBGraph. Since they are constructed from a
ConceptBase database, they are a (materialized) view on the database.
The view consists of the nodes and links displayed in the window, their positions,
their graphical types, the hostname, portnumber and module from which the graph was created,
the size of the window, its background color and image, and the window's zoom factor. 
You can thus save the current state of your graphical view in the {\em gel\/} file
and load it in a subsequent session with CBGraph to continue editing it,
much like with a drawing tool. 

The graph file stores serialized Java objects in the following sequence
\begin{verbatim}
   String title of the internal frame
   Dimension size of the graph editor
   Dimension size of the internal window
   Dimension size of the drawing area of the internal window
   Integer number of nodes (incl. edge objects)
     { node1
       Rectangle bound of node1
       node2
       Rectangle bound of node2
       ...
     }
   Integer number of edges
     { String source node of edge 1
       String node object on the edge 1 
       String destination of edge 1
       String source node of edge 2
       String node object on the edge 2 
       String destination of edge 2
       ...
     }
   Color background color
   Float zoomfactor
   String hostname
   String port number
   String module context as absolute path
   String long title of the graphical palette 
   String name of the graphical palette
   Integer saveflag (bit0: HAS_IMAGE, bit1: HAS_SOURCE, bit2: HAS_PARAMS)
   IF HAS_PARAMS 
     { String[] params
     }
   IF HAS_SOURCE
     { String saved modules e.g. "mod1-mod2"
       { String source of module 1
         String source of module 2
         ...
       }
     }
   IF HAS_IMAGE 
     { PNG image of the background image
     } 
\end{verbatim}

Since edges are also objects in Telos, the nodes stored in the graph file 
include the edge node that represents the edge itself.
The graph file stores complex information about the nodes including
the graphical types of the nodes, their dimension and location.
By default, the graph files stores the Telos module sources needed to manipulate
the graph. The graph file stores the sources of all modules that are on the path
from the root module {\tt System} to the current module (the module that is active when
the graph file is saved). The module {\tt System} is not saved since it is typically
not updated. Note that ConceptBase database can include a tree of module
(see section \ref{sec:nested_modules}).
Hence, a graph file does in general not store the Telos sources of the complete
database. If you create graph files for all leave modules of a database, then
the combination of the graph files is completely containing the database as sources models.
The extraction of Telos sources uses the builtin query {\tt listModule}
(see section \ref{sec:module_content}). This is in most cases a faithful listing. However,
there are rare cases when the extracted cannot be told to a CBserver.
For example, if a module contains deductive rules that are essential for
satisfying integrity constraints for objects defined in the same module, then
a single TELL operation could fail bacause ConceptBase requires the deductive rules to
be compiled. See section \ref{sec:listmod} for more details.

If you specified command line parameters like "+r", "-r", "+rw", or "-rw" at the start of CBGraph,
then these parameters are stored in the GEL file.

The background image is not stored as a serializable Java object but as a PNG
image using the ImageIO class of Java. It is always stored as the last element since
the input routines shall read it until the end of the file.
Note that some strings can be just null.


\section{An example session with ConceptBase}
\label{sec:session}


In this section we demonstrate the usage of the {\em ConceptBase.cc User Interface}, by involving an example model. It consists of a few classes
including {\tt Employee, Department, Manager}. The class {\tt
Employee} has the attributes {\tt name, salary, dept, and boss}. In
order to create an instance of {\tt Employee} one may specify the
attributes {\tt salary, name,  and dept}. The attribute {\tt boss}
will be computed by the system using the {\tt bossrule}. There is also a
constraint which must be satisfied by all instances of the class {\tt
Employee} which specifies that no employee may earn more money than its
boss. The Telos notation for this model is given in Appendix
\ref{sec:employee-model}.


\subsection{Starting ConceptBase}


To start a {\em ConceptBase} session, we use two terminal windows, one for the {\em ConceptBase.cc}
server and one for the usage environment. We start the {\em ConceptBase.cc} server
by typing the command 
\begin{verbatim}
   cbserver -port 4001 -d test  
\end{verbatim}
in a terminal window of, let us say machine alpha%
\footnote{A full list of all parameters is described in section \ref{cha:cbserver}.
Note that the script {\tt CBserver} must be in the search path. It is available in
the subdirectory {\tt bin} of your ConceptBase.cc installation directory.}%
. The parameter -port sets the port number under which the CBserver communicates to
clients and the parameter -d specifies the name of the directory into which the 
CBserver internally stores objects persistently.
Then, we start the usage environment with
the command {\tt cbiva} in the other window. 

It is also possible to start the CBserver from the user interface.
To do so, choose ``Start CBserver'' from the ``File Menu'' of CBIva (see
section \ref{sec:menubar}) and specify the parameters in the dialog which
will be shown (see figure \ref{fig:startserver}).
The option {\tt Source Mode}  controls whether the CBserver accessing
the database via the {\tt -d} parameter (database maintained in binary files),
or via the {\tt -db} parameter (database maintained both in binary files and
in source files). See section \ref{cha:cbserver} for more details.
Once the information has been
entered via the OK button, the server process will be started and its output will be captured
in a window. This output window provides also a button stop the server.
If you started the server this way then you can skip the next section,
as the user interface will be connected to the server automatically.

\cbfigure{startserver}{0.9\textwidth}{Start CBserver dialog and CBserver output window}



\subsection{Connecting CBIva to another CBserver}

By default, CBIva will automatically connect to a local or public CBserver when started.
If you want to start a CBserver with dedicated parameters from CBNIva, then first 
select {\bf File/Disconnect} and then
establish a new connection between the {\em ConceptBase.cc}  server and the user
client CBIva. This is done by choosing the option {\bf Connect}
from the File menu of CBIva. An interaction window
appears (see Figure \ref{fig:connectserver}) querying for the host name
and the port number of the server (i.e. the number we have specified within
the command {\tt cbserver -port 4001 -d test}).

\cbfigure{connectserver}{3cm}{The connect-to-server dialog}

\subsection{Loading objects from external files}

The objects manipulated by {\em ConceptBase.cc} are persistently stored in a
collection of external files, which reside in a directory called {\bf
application} or {\bf database}%
\footnote{Historically, we used the terms 'application' or 'object base' instead 'database'.
We now believe that 'database' is a much better term.}%
. The actual directory name of the
database is supplied as the -d parameter of the command {\tt
CBserver}.

The -u parameter of the {\tt CBserver} specifies whether updates are made
persistent or are just kept in system memory temporarily. Use {\tt -u
persistent} for a update persistence or {\tt -u nonpersistent} for a
non persistent update mode \footnote{In nonpersistent update mode, the database
is actually copied to a temporary directory. This copy will be removed
when you shutdown the server.}.

The database can be modified interactively using the editor commands
TELL/UNTELL.
Another way of extending ConceptBase.cc databases is to load Telos objects (expressed
in frame syntax) stored in plain text files with the extension  *.sml. Call the menu
item {\bf Load Model} from the {\bf File Menu} to add these objects
to the database. In our example the database (directory)
Employee was built interactively and can be found together with files
containing the frames constituting the example in the directory

\begin{center}	{\em CB\_HOME}{\tt /examples/QUERIES}
\end{center}

\noindent where you have to replace {\em CB\_HOME} with the {\em ConceptBase}  installation directory.  The
following files contain the objects of the Employee example expressed in
frame syntax:  {\tt Employee\-\_Classes.sml, Employee\-\_Instances.sml,
 Employee\-\_Queries.sml}.  An alternative to interactively building a
database is to start the server with an empty database ({\tt -d
}$\langle ${\tt newfile}$\rangle $) and  then add the objects in these files
by using {\bf Load Model}. Note, that the *.sml extension may be omitted.
 During the load operation of external models, {\em ConceptBase}
checks for syntactical and semantical correctness and reports all errors to
the history window as it is done when updating the object base interactively using the
editor. This protocol field collects all operations and errors reported since the beginning of
the session.



\subsection{Displaying objects}


To display all instances of an object, e.g. the class {\tt Employee}, we invoke
the Display Instance facility  by
selecting the item {\bf Display Instances} from the menu bar. In the
interaction window we specify {\tt Employee} as object name. The
instances of the class {\tt Employee} are then displayed (see Figure \ref{fig:session_display}).


\cbfigure{session_display}{7cm}{Display of {\tt Employee} instances}

After selecting a displayed instance we can load the frame representation
of an instance to the Telos Editor or display further instances.

\subsection{Browsing objects}

The {\em Graph Editor\/} is the preferred tool for browsing the objects managed
by a CBserver.
CBGraph is started by using the menu item ``Graph Editor'' from the ``Browse''
menu of CBIva. Select {\tt Employee} as the initial object to be shown in CBGraph.
After starting CBGraph, it will open an internal frame, connect
it with the current server, and load the {\tt Employee} object.

The {\em CBGraph Editor} (described in detail in section \ref{sec:graphed}) allows you to
display arbitrary objects from the current onceptBase server.
Then, we select the {\tt Employee} object and choose the {\bf Sub classes} option from the context menu available via
the right mouse-button. We choose to only display explicit subclasses from the submenu and
select the {\tt Manager} object.
The displayed graph is now expanded (figure \ref{fig:graph_empl}).


\cbfigure{graph_empl}{0.9\textwidth}{The resulting graph after expanding the node {\tt Employee} with subclasses}

Now we expand the node {\tt Manager}, a subclass of {\tt Employee}, by choosing the
menu item {\bf Instances} from the popup menu for Manager. We select the menu item
``Show all'' to display all instances of {\tt Manager}.
The resulting graph is shown in figure \ref{fig:graph_man}.

\cbfigure{graph_man}{0.9\textwidth}{The resulting graph after expanding with the instances of {\tt Manager}}

Note that different object types are represented by different graphical objects.
The instances of {\tt Manager} are
shown only as grey rectangles, because they are normal individual objects.
The nodes {\tt Manager, Salesman, Employee} etc. are shown as ovals,
since these nodes are instances of the system class {\tt SimpleClass} (see
 for a full description of  graphical
object semantics: Appendix \ref{cha:graph-typen}).


One can move nodes and links by selecting a node or a link and then holding
down the left mouse button while moving the cursor to a different
position. When the button is released the selected object will be located
at the current position and the related links are redisplayed.
Selection and movement of multiple nodes and links is also possible.
Nodes and linkes are selected by clicking on its label, e.g. {\tt Manager}
in figure \ref{fig:graph_man}. Some links like the blue specialization
link in the figure have no label. Then one can select the link by clicking on the
small square dot in the middle of the link. This square dot is by default invisible.
It becomes visible when you click on any other node or link label in the graph,
e.g. on {\tt Manager}. 

We can further experiment with CBGraph by showing
the classes and attributes of Employee. The classes of Employee are
shown by selecting ``Instance of'' from the popup menu.
Attributes of an object can be shown by selecting ``Outgoing attributes''
from the menu. The next submenu will show all attribute classes that apply
for the current object. In our example, {\tt Employee} is an instance of
{\tt Class}. Therefore, it has the attribute classes {\tt constraint}, {\tt rule},
and {\tt mrule} (see figure \ref{fig:graph_super}). The attribute class {\tt Attribute} applies to all objects
as in Telos any kind of object can have an attribute. Furthermore, all attributes
of an object are member of the attribute class {\tt Attribute}. As we want
to see all attributes, we select this attribute class and select ``Show all''
from the next submenu. All attributes and their values will be shown
in CBGraph.

\cbfigure{graph_super}{0.9\textwidth}{The graph after expanding it with the classes and the attributes of the class {\tt Employee}}

CBGraph can also show implicit relationships between objects,
e.g. relationships deduced by rules or the Telos axioms.
For example, if we select the object {\tt John} and select ``Instance of''
from the popup menu, we can display the implicit classes of {\tt John}
by selecting ``All'' from the next submenu. As {\tt John} is an instance of
{\tt Manager} and {\tt Manager} is a subclass of {\tt Employee}, {\tt John} is also an instance
of {\tt Employee}. As there is no explicit object \verb+John->Employee+,
the instantiation link between {\tt John} and {\tt Employee} will be represented as
an implicit link, i.e. a dashed line (see figure \ref{fig:graph_impl}).

The same applies also to attribute links. For example, the employee {\tt Herbert}
has an implicit {\tt boss}-attribute to {\tt Phil}. This can be shown by selecting
``Outgoing attributes'' $\rightarrow$ ``boss'' $\rightarrow$ ``All'' $\rightarrow$ ``Phil'' from the
popup menu. Note, that the submenu ``All'' for the attribute class {\tt Attribute}
will be always empty as only explicit attributes can be displayed in this category.

\cbfigure{graph_impl}{0.9\textwidth}{The graph showing implicit instantiation and attribute links}

\subsection{Editing Telos objects}


\subsubsection{Editing Telos objects with the Telos editor}

Before we are able to edit a Telos object, we have to load its frame
representation in to the Telos Editor field first. For loading a Telos
object to the Editor field, we choose the {\bf Telos Editor} Button from
either the Display Queries oder Display Instances Browsing facilities or the
{Load Frame} button from the {\em ConceptBase\/Workbench}\ window (see
Figure \ref{fig:telos_ed}).


\cbfigure{telos_ed}{0.8\textwidth}{The TelosEditor Field with the object {\tt Employee}}

Now we add an additional attribute, e.g. {\tt education}, to the class {\tt Employee}
(for the description of the Telos syntax see Appendix \ref{cha:syntax}). We have added
the line {\tt education : String} as shown in figure \ref{fig:telos_ed_error}.
To demonstrate error reports
from the {\em ConceptBase.cc\/} user interface and how to correct them, we have made mistakes in the syntax
notation of the added attribute.

\cbfigure{telos_ed_error}{13cm}{Trying to add an attribute to the class {\tt Employee} with the resulting error report}

By clicking the left mouse button on the {\bf Tell} icon, the content of the
editor is told to the {\em ConceptBase}  server. Syntactical and semantical
correctness is checked and the detected errors are reported to the Protocol field.
The report resulting from our mistakes by specifying
the new attribute is also shown in figure \ref{fig:telos_ed_error}.
Note, that this syntax error would have been already detected at the client side
without interaction with the server if we would have enabled the option
``Pre-parse Telos Frames'' in the options menu.

We correct the error by adding a semicolon to the previous line
and choose the {\bf Tell} symbol again. This time,
since there are no further mistakes, the additional attribute is added to
the class {\tt Employee}.

Now we can choose again the item {\bf Outgoing Attributes} from the popup of the
{\em Graph Editor} window for the node Employee. If we select ``Show all'' attributes
of the attribute class ``Attribute'' the new attribute will NOT be shown.
CBGraph uses an internal cache which will only be updated on request.
Therefore, we select the object {\tt Employee} and select ``Validate and
update selected objects'' from the menu ``Current connection''. The cache
for the object {\tt Employee} will be emptied. Now, showing all attributes
should add the new attribute {\tt education} to the graph.

\subsubsection{Editing Telos objects using CBGraph}

Telos objects can also be edited graphically in CBGraph. In our
example, we want to add another attribute named {\tt address} to the class {\tt Employee}.
The attribute destination of this attribute should be a new class called {\tt Address}.


First, we select the object Employee and click on the ``Create Attribute'' button in
the tool bar or select ``Add Attribute'' operation from the tool bar. As we have
selected the Employee object, it should be already inserted as source of
the attribute. We have to type the label (address) and the destination
of the attribute in the text fields (see figure \ref{fig:graph_add_attr}).
As this attribute does not belong to a specific attribute category (it is
just an attribute), we do not have to specify an attribute class.

By clicking on the Ok button, CBGraph will create a new object
for {\tt Address} represented by the default graphical type (a gray box).
Then, it will create the attribute link from {\tt Employee} to {\tt Address}
with the label {\tt address}. The result is shown in figure \ref{fig:graph_add_inst}.

Now, we want to declare {\tt Address} as an instance of {\tt Class}. Therefore, we
select {\tt Address}, hold down the Shift-key and select {\tt Class}. Both objects should
be selected now. We click on the ``Create Instantiation'' button and a dialog as shown
in figure \ref{fig:graph_add_inst} should appear.

\cbfigure{graph_add_attr}{0.8\textwidth}{Adding an attribute to {\tt Employee}}

\cbfigure{graph_add_inst}{0.8\textwidth}{Adding an instantiation link between {\tt Address} and {\tt Class}}

As we have selected the objects in the correct way, the dialog already specifies the
object we want to create ({\tt Address} in {\tt Class}) and we can click directly on Ok.
The new instantiation link will be added to the graph.

\cbfigure{graph_commit}{0.8\textwidth}{The resulting graph after commit}

Now, we are satisfied with our changes and want to commit them in the server. So far, the
changes have been stored in an internal buffer of CBGraph and have not been
sent to the server. We click on the ``Commit'' button in the upper right corner.
CBGraph generates now Telos frames for the added objects and sends them
to the server. If we did not make an error, all changes should be consistent
and accepted by the server. This is shown by the appearing message box ``Changes committed''.
If an error occurs, an error message will displayed instead. The graphical editing is
an alternative to the textual editing via the Telos Editor. It is appropriate for incremental
changes to a model. Larger changes should better be made via the Telos Editor or
even to an external text file that is loaded via the {\em File / Load Model\/} facility
of CBIva.
If the changes were successfully told to the CBserver, CBGraph
reloads the information of every visible object. In particular, the graphical
types of the objects will be updated. As the object {\tt Address} is declared
as instance of {\tt Class}, it will get the graphical type of a class, i.e.
a turquoise box. The result is shown in figure \ref{fig:graph_commit}. The object
{\tt Employee} is shown in the detailed component view, in this case the
frame representation of the object is shown. As you can see, the attribute
{\tt address} is now also visible in the frame representation.



\subsection{Using the query facility}


Lets assume that we need to ask the server for all {\tt Employees} working for {\tt Angus}. 
We open a new Telos Editor (see menu item Browse).
Then, we define a new query class ({\tt AngusEmployees}) as follows:

\begin{verbatim}
AngusEmployees in QueryClass isA Employee with
constraint
    c: $ (this boss Angus) $
end
\end{verbatim}

We can tell this query, so that it is stored in ConceptBase.cc and we can reuse it later,
or we can just {\em ask} the query, i.e.\ the query will told temporarily and evaluated.
If we ask the query, the answer is displayed in the Telos Editor field as well as in
the history window. Figure \ref{fig:query_answer} shows the
CBIva with the query class and the answer.

\cbfigure{query_answer}{0.8\textwidth}{Query class and its answer}

We can also execute this query from CBGraph. From the menu ``Current connection''
we select ``Query to server''. A dialog we ask for the name of query class. If we have
told the example query, we can type {\tt AngusEmployees} in the text field and
hit on the ``Submit Query''. The query will be evaluated and the objects in the result
will be shown in the list box. We can select the objects which should be added to the graph
(multiple selection with the Shift-key is possible) and click on the ``Show objects'' button.
The selected objects will be added to the graph, however with no connection to existing objects.


\section{Configuration file}
\label{sec:config}

The configuration options are stored in a file ``.CBjavaInterface'' in
the home directory of the user. The settings are stored automatically on exit.
You can edit the file manually with a normal text editor. It contains
name-value pairs in the format \verb+variable=value+. All variables
can also be modified CBIva via the "Options" menu.

\begin{itemize}
\item Variables related to CBIva
\begin{description}
\item[PathForLoadModel:] Path used by the load model dialog (contains the most recent directory selected in a dialog).
\item[RecentConnections:] Comma-separated list of recent connections in the format host/port, applies also to CBGraph.
\item[PreParseTelosFrames:] Frames are be parsed on client-side before sent to the server (true/false).
\item[UseQueryResultWindow:] Use the query result window to display results of a query (true/false).
\item[ConnectionTimeout:] Number of milliseconds the interface waits for a response of the server
\item[LPICall:] Enable LPI-Call (internal use only).
\item[ShowLineNumbers:] If set to 'true', the Telos editor of CBIva shall display line numbers.
\item[CBIvaSmallfont:] The font size of the TelosEditor text and log area as regular (small) font, default "12f".
\item[CBIvaLargefont:] The font size of the TelosEditor text and log area as large font, default "18f".
\item[PublicCBserver:] Either 'none' (=disabled) or the domain name of a computer that hosts a publicly accessible
ConceptBase server. A port number can be appended as well like in 'cbserver.acme.com/4002'. Default port number is 4001.
The variable PublicCBserver also applies to CBGraph and CBShell. See section \ref{sec:pubcbserver} for more details. If the value is different
from 'none', then CBIva shall attempt to connect to the public CBserver at startup.
\end{description}


\item Variables related to CBGraph
\begin{description}
\item[PathForLayout:] Path used by the dialog to store and load graphs (contains the most recent directory selected in a dialog).
\item[ComponentView:] Default view for the detailed representation in CBGraph (might be ``\verb+frame+'' or ``\verb+tree+'').
\item[SaveLayoutWithGraphType:] Layouts of CBGraph are stored with all information about graphical types (true/false).
\item[InvalidObjectsMethod:] Specifies whether objects that have been identified as invalid should be marked or deleted (mark/delete).
\item[DiagramDesktopBackgroundColor:] Background color of the desktop of CBGraph (comma-separated representation of an RGB-value).
\item[DebugLevel:] Level for debug messages. Possible values are SEVERE, WARNING, INFO, CONFIG, FINE, FINER, FINEST
(according to the java.util.logging package). Default is WARNING.
\item[ModuleSeparator:] Can be either '-' or '/'. 
\item[NodeLevelAware:] Enables or disables the special behavior of nodes with negative level in CBGraph. See also
section \ref{sec:nodelevel}. Default is "true".
\end{description}
\end{itemize}


